\begin{table}[htbp]
\centering
\caption{Distribuição dos potenciais pares por nota final do comparador}
\label{tab:score-bands}
\begin{tabular}{lrrr}
\toprule
Nota final do comparador & Total & Pares verdadeiros & \% pares verdadeiros \\
\midrule
    0-3 & 5.336 & 0 & 0,0\% \\
    3-5 & 34.645 & 3 & $<$0,1\% \\
    5-6 & 14.263 & 17 & 0,1\% \\
    6-7 & 6.175 & 57 & 0,9\% \\
    7-8 & 1.080 & 43 & 4,0\% \\
    8-9 & 102 & 34 & 33,3\% \\
    9-10 & 44 & 43 & 97,7\% \\
    $\geq$10 & 51 & 50 & 98,0\% \\
\bottomrule
\end{tabular}

\smallskip
{\footnotesize\textit{Nota}: a tabela utiliza a coluna \texttt{nota final} do comparador, e não o escore probabilístico do OpenRecLink. As faixas são fechadas à esquerda e abertas à direita; a última inclui notas a partir de 10. Fonte: extrato pseudonimizado, coluna \texttt{nota final}, e \texttt{scripts/generate\_missing\_figures.py}.}
\par\footnotesize Total: 61.696 potenciais pares, 247 pares verdadeiros. Fonte agregada: \path{data/tab_score_bands_proveniencia.csv}. Par verdadeiro significa PAR igual a 1 ou 2.
\end{table}
